\begin{lemma}[Scheide maximal-tree package (d)]
Under the fixed-pair maximal-tree hypotheses, every Tashkinov tree $S$ for
$(e_0,\varphi)$ is an initial segment of a Tashkinov tree $\widetilde T$
with $V(\widetilde T)=V(T)$.
\end{lemma}
\begin{proof}
Write
\[
 T=(x_0,f_1,x_1,\ldots,f_m,x_m),\qquad X=V(T).
\]
By \noderef{ScheideMaximalContainsEveryTree}, $V(S)\subseteq X$.  We
construct a sequence of enlargements, beginning with $C_0=S$, and maintain
the following three assertions about the current sequence $C$: it is a
Tashkinov tree, $S$ is an ordered initial segment of $C$, and
$V(C)\subseteq X$.  The meanings of the two vertex sets here are those of
\noderef{TreeSequenceVertexSet}.  These assertions hold for $C_0$: the
first is the given Tashkinov-tree hypothesis on $S$, the second follows
from \noderef{TreeSequenceInitialSegment} because $S$ has its own length
and agrees with itself entry by entry, and the third is the inclusion just
obtained.

Consider one stage, and display the current sequence as
\[
 C=(z_0,g_1,z_1,\ldots,g_q,z_q).
\]
The definition \noderef{TreeSequence} gives $q\geq1$, $g_1=e_0$, and says
that the ends of $g_1$ are $z_1$ and $z_0$.  Applied to $T$, the same
definition gives $f_1=e_0$ and says that the ends of $f_1$ are $x_1$ and
$x_0$.  Consequently the unordered endpoint pair of $e_0$ is both
$\{z_0,z_1\}$ and $\{x_0,x_1\}$.  In particular, each of $x_0,x_1$ occurs
among the current vertices.  If $V(C)\ne X$, let $j$ be the least index for
which $x_j\notin V(C)$.  The preceding observation gives $j\geq2$, and
minimality gives
\[
 x_h\in V(C)\quad\hbox{for every }h<j. \tag{1}
\]

The attachment clause of \noderef{TreeSequence} for the original edge
$f_j$ gives an index $i<j$ such that the ends of $f_j$ are $x_j$ and
$x_i$.  By (1), $x_i$ has some current index $r$, say $z_r=x_i$; this is
exactly the membership description in \noderef{TreeSequenceVertexSet}.
The edge identity $f_j$ is not one of $g_1,\ldots,g_q$.  Indeed, if
$f_j=g_\ell$, the attachment clause for the already present identity
$g_\ell$ would say that both ends of this identity are current vertices.
Its endpoint pair is also $\{x_j,x_i\}$, so $x_j$ would equal one of those
two current endpoints.  This contradicts the choice of $j$.

Append $f_j,x_j$ at the end of $C$, and call the resulting sequence $C^+$.
We check all clauses of \noderef{TreeSequence}.  Its length is $q+1$, still
positive.  Its old vertices remain pairwise distinct and the new vertex is
distinct from all of them because $x_j\notin V(C)$, so its vertex map is
injective.  Its old edge identities remain pairwise distinct and $f_j$ is
not among them by the preceding paragraph, so its edge map is injective.
The first edge and every old attachment are unchanged.  Finally, the new
edge has ends $x_j$ and $x_i=z_r$, where the current index satisfies
$r\leq q<q+1$; hence it attaches the new last vertex to an earlier vertex.
Thus $C^+$ is a tree-sequence.  Since it was formed only by appending its
last edge and vertex, $C$ is an ordered initial segment of $C^+$ in the
sense of \noderef{TreeSequenceInitialSegment}.

It remains to check the Tashkinov condition.  By
\noderef{FixedPairMaximalTashkinovTree}, $T$ is itself a Tashkinov tree.
For each old edge of $C^+$ whose index is at least two, use the witness it
had in the current Tashkinov tree $C$.  The edge and every earlier vertex
retain the same indices in $C^+$ because $C$ is its ordered initial segment,
so this is still a valid witness.  For the new edge, $j\geq2$ permits us to
apply the defining witness clause of \noderef{TashkinovTree} to $f_j$ in
$T$.  It supplies an index $h<j$ and a color $\alpha$ such that
$\varphi(f_j)=\alpha$ and $\alpha$ is missing at $x_h$.  By (1), $x_h$ has
some current index $s\leq q$.  In $C^+$ that same vertex has index $s$, and
$s<q+1$, the index of the newly appended vertex.  Hence it is the required
earlier missing-color witness for $f_j$.  The old and new cases together
are precisely the condition in \noderef{TashkinovTree}, so $C^+$ is a
Tashkinov tree.  Moreover, no entry of $S$ was changed: its entries agree
with those of $C$, and those of $C$ agree with the initial entries of
$C^+$.  Thus $S$ remains an ordered initial segment.  Finally,
\[
 V(C^+)=V(C)\cup\{x_j\}\subseteq X, \tag{2}
\]
where the equality follows directly from the entry-wise definition in
\noderef{TreeSequenceVertexSet}.

To iterate this step rigorously, use induction on the finite number
\[
 \mu(C)=|X\setminus V(C)|.
\]
If $\mu(C)=0$, the finite set $X\setminus V(C)$ has cardinality zero and
is therefore empty, so no enlargement is needed.  Otherwise the least
absent vertex $x_j$ exists, and the construction above applies.  Since
$x_j$ lies in $X$ but not in $V(C)$, while (2) shows that it is the only
new vertex, $X\setminus V(C^+)$ is obtained from $X\setminus V(C)$ by
deleting exactly $x_j$.  Thus $\mu(C^+)=\mu(C)-1$.  The three maintained
assertions were verified above, so the induction continues and terminates
at a Tashkinov tree $\widetilde T$ with
$X\setminus V(\widetilde T)=\varnothing$.

We now verify the required vertex-set equality by both inclusions.  If
$v\in V(\widetilde T)$, then either $v$ was already a vertex of $S$, in
which case $v\in X$ by the initial inclusion, or $v$ was appended as one of
the selected vertices $x_j$ of $T$, again giving $v\in X$.  Hence
$V(\widetilde T)\subseteq X$.  Conversely, if $v\in X$ but
$v\notin V(\widetilde T)$, then $v\in X\setminus V(\widetilde T)$,
contrary to its emptiness.  Hence $X\subseteq V(\widetilde T)$, and so
$V(\widetilde T)=V(T)$.  At every induction step entries were appended only
at the end.  Therefore the final sequence has at least the length of $S$
and agrees with every vertex and edge entry of $S$ through that length,
which is exactly
$\noderef{TreeSequenceInitialSegment}$ for $S$ and $\widetilde T$.
\end{proof}
